% ==============================================================================
% 08_incident_response.tex — UE SIS589
% ==============================================================================
\chapter{Réponse aux Incidents de Sécurité}
\label{chap:ir}

\epigraph{\itshape « La préparation n'est pas une dépense~;
  c'est une assurance contre le chaos. »}
         {--- NIST SP~800-61}

\begin{boxobjectifs}
\begin{enumerate}
  \item Dérouler le cycle de réponse aux incidents selon ISO~27035 et NIST~SP~800-61.
  \item Constituer et opérer un CSIRT~; définir ses rôles et ses interfaces.
  \item Rédiger un playbook opérationnel pour un type d'incident courant.
  \item Appliquer les techniques de confinement, d'éradication et de
        rétablissement.
  \item Mapper les phases de réponse aux tactiques MITRE ATT\&CK.
  \item Rédiger un rapport d'incident complet.
\end{enumerate}
\end{boxobjectifs}

% ==============================================================================
\section{Le cycle de réponse aux incidents}
\label{sec:ir-cycle}
% ==============================================================================

\subsection{ISO~27035 et NIST~SP~800-61}

Les deux référentiels s'articulent autour des mêmes phases fondamentales~:

\begin{table}[H]
\centering\small
\caption{Comparaison ISO~27035 et NIST~SP~800-61}
\label{tab:ir-cycle}
\rowcolors{2}{TableauPair}{TableauImpair}
\begin{tabular}{p{2.5cm}p{5cm}p{5.5cm}}
\toprule
\tableauHeader{Phase} & \tableauHeader{ISO~27035} &
\tableauHeader{NIST SP~800-61} \\
\midrule
1 & Planification et préparation &
Préparation (\textit{Preparation}) \\
2 & Détection et signalement &
Détection et analyse (\textit{Detection \& Analysis}) \\
3 & Évaluation et décision &
(inclus dans Détection et analyse) \\
4 & Réponse (confinement, éradication, rétablissement) &
Confinement, éradication, rétablissement \\
5 & Leçons tirées &
Activités post-incident \\
\bottomrule
\end{tabular}
\end{table}

\subsection{Phase~1~: Préparation}

La préparation est la phase la plus déterminante~: la qualité de la
réponse en situation de crise est proportionnelle à la qualité de la
préparation en situation calme.

Éléments de préparation indispensables~:
\begin{itemize}
  \item \textbf{Politique de gestion des incidents~:} définit ce qu'est
        un incident, les niveaux de gravité, les procédures d'escalade,
        les obligations légales de notification.
  \item \textbf{CSIRT constitué~:} équipe identifiée, formée, disponible
        en astreinte. Contacts d'urgence actualisés.
  \item \textbf{Playbooks rédigés~:} procédures pour chaque type d'incident
        probable (ransomware, phishing, fuite de données, DDoS).
  \item \textbf{Outils préparés~:} station forensique SIFT, supports
        d'acquisition vierges, bloqueurs en écriture, accès SIEM.
  \item \textbf{Exercices réguliers~:} simulations Tabletop, red team.
\end{itemize}

\subsection{Phase~2~: Détection et analyse}

\subsubsection{Sources de détection}

\begin{table}[H]
\centering\small
\caption{Sources de détection des incidents}
\label{tab:detection-sources}
\rowcolors{2}{TableauPair}{TableauImpair}
\begin{tabular}{p{4cm}p{4cm}p{5.5cm}}
\toprule
\tableauHeader{Source} & \tableauHeader{Type} &
\tableauHeader{Fiabilité / Délai} \\
\midrule
Alerte SIEM & Automatique & Rapide, taux de FP variable \\
Alerte EDR / AV & Automatique & Fiable pour les malwares connus \\
Signalement utilisateur & Manuel & Tardif mais indicateur de gravité \\
Banque / partenaire & Externe & Fiable, souvent tardif \\
CERT national / CERT-MU & Externe & Haute fiabilité, délai variable \\
Threat hunting & Proactif & Excellent~; nécessite N3 disponible \\
\bottomrule
\end{tabular}
\end{table}

\subsubsection{Triage et qualification}

L'analyste N1 répond à trois questions~:
\begin{enumerate}
  \item \textbf{Est-ce un vrai positif~?} (ou un faux positif)
  \item \textbf{Quelle est la gravité réelle~?} (impact technique + impact métier)
  \item \textbf{Quelle est l'urgence~?} (contenir maintenant ou analyser d'abord)
\end{enumerate}

\begin{table}[H]
\centering\small
\caption{Grille de classification des incidents LiZbiZ-SIS}
\label{tab:gravite-incidents}
\rowcolors{2}{TableauPair}{TableauImpair}
\begin{tabular}{p{1.8cm}p{4cm}p{4.5cm}p{3.5cm}}
\toprule
\tableauHeader{Gravité} & \tableauHeader{Critères} &
\tableauHeader{Exemples} & \tableauHeader{Délai de réponse} \\
\midrule
\gravite{critique} &
Impact immédiat sur la continuité ou données très sensibles &
Ransomware, exfiltration BDD clients, compromission AD &
$<$~1~heure, astreinte activée \\
\gravite{haute} &
Impact significatif, propagation probable &
Brute force réussi, malware détecté non confiné &
$<$~4~heures \\
\gravite{moyenne} &
Impact limité, propagation improbable &
Phishing réussi sur un poste, scan interne &
$<$~24~heures \\
\gravite{basse} &
Événement de sécurité sans impact confirmé &
Tentative de phishing bloquée, scan externe &
$<$~72~heures \\
\bottomrule
\end{tabular}
\end{table}

% ==============================================================================
\section{Le CSIRT}
\label{sec:csirt}
% ==============================================================================

\begin{boxdef}
\textbf{CSIRT (\textit{Computer Security Incident Response Team})~:}
Équipe dédiée à la réception, au tri, à l'analyse, à la réponse et à la
coordination lors d'incidents de sécurité. Point de contact unique.
Un CSIRT interne est différent d'un CERT national (externe, sectoriel).
\end{boxdef}

\begin{table}[H]
\centering\small
\caption{Rôles dans un CSIRT}
\label{tab:csirt-roles}
\rowcolors{2}{TableauPair}{TableauImpair}
\begin{tabular}{p{3.5cm}p{5.5cm}p{4.5cm}}
\toprule
\tableauHeader{Rôle} & \tableauHeader{Responsabilités} &
\tableauHeader{Interfaces} \\
\midrule
Incident Manager & Coordonne l'ensemble de la réponse~; tient le log
                    d'incident~; décide des escalades &
DG, RSSI, équipes \\
Analyste technique & Investigation forensique, confinement technique,
                     éradication &
SOC N2/N3, équipes IT \\
Responsable communication & Communication interne (direction, métiers)
                             et externe (clients, presse, autorités) &
DG, communication, CNIL \\
Coordinateur légal & Évalue les obligations de notification~;
                      interface avec les juristes et le parquet &
Juriste, CNIL, ANTIC \\
\bottomrule
\end{tabular}
\end{table}

\begin{boxlizbiz}
\textbf{LiZbiZ-SIS~: CSIRT minimal}

LiZbiZ-SIS n'a pas de CSIRT formel. Suite à l'incident du 12/04/2026,
la DSI a constitué une cellule de crise ad hoc~: 2~analystes SOC
(N1 + N2), 1~ingénieur réseau, le RSSI (à temps partiel),
1~représentant de la direction. Une convention avec un MSSP a été
signée pour l'assistance forensique de niveau N3. Un numéro d'astreinte
a été créé. Prochaine étape~: rédiger les playbooks ransomware et
fuite de données.
\end{boxlizbiz}

% ==============================================================================
\section{Playbooks et Runbooks}
\label{sec:playbooks}
% ==============================================================================

\begin{boxdef}
\textbf{Playbook~:} Document décrivant la procédure logique de réponse
à un type d'incident donné. Répond à \og{}Quoi faire~?\fg{} en termes de
décisions, escalades et communications. Indépendant de la technologie.
\end{boxdef}

\begin{boxdef}
\textbf{Runbook~:} Document technique détaillant les commandes exactes
et les actions précises à exécuter pour chaque étape du playbook.
Répond à \og{}Comment faire exactement~?\fg{}. Technologie-spécifique.
\end{boxdef}

\begin{boxplaybook}
\textbf{Playbook~: Compromission de compte (accès non autorisé SSH)}

\textbf{Déclencheur~:} Alerte SIEM règle~100200 (brute force SSH réussi)
ou signalement utilisateur d'activité inhabituelle.

\textbf{Étape~1~--- Qualification (N1, $<$~15~min)~:}
\begin{itemize}
  \item Vérifier l'alerte SIEM~: l'IP source est-elle dans la liste blanche~?
  \item Identifier l'utilisateur concerné et son profil (admin / standard).
  \item Si faux positif confirmé~: fermer le ticket, documenter.
  \item Si incident confirmé~: escalader N2, activer le log d'incident.
\end{itemize}

\textbf{Étape~2~--- Confinement immédiat (N2, $<$~30~min)~:}
\begin{itemize}
  \item Désactiver le compte compromis dans l'AD.
  \item Bloquer l'IP source au firewall pfSense.
  \item Isoler la machine cible dans un VLAN de quarantaine.
  \item Capturer l'état de la mémoire (avml / WinPmem) avant tout arrêt.
\end{itemize}

\textbf{Étape~3~--- Investigation (N2/N3, $<$~4~h)~:}
\begin{itemize}
  \item Analyser auth.log~: depuis quand~? Quelle IP~? Quel user~?
  \item Analyser la timeline Sysmon~: commandes exécutées, fichiers créés,
        connexions sortantes.
  \item Vérifier les mouvements latéraux depuis cette machine.
  \item Identifier les persistances (cron, Run keys, nouveaux comptes).
\end{itemize}

\textbf{Étape~4~--- Éradication~:}
\begin{itemize}
  \item Supprimer les artefacts malveillants identifiés.
  \item Réinitialiser les credentials compromis.
  \item Patcher la vulnérabilité exploitée.
\end{itemize}

\textbf{Étape~5~--- Rétablissement~:}
\begin{itemize}
  \item Remettre la machine en production (restauration si nécessaire).
  \item Renforcer la surveillance (nouvelle règle SIEM ciblée).
\end{itemize}

\textbf{Étape~6~--- Notification~:}
\begin{itemize}
  \item Notifier la direction~: impact, mesures prises.
  \item Évaluer l'obligation de notifier l'ANTIC.
  \item RETEX dans les 72~heures.
\end{itemize}
\end{boxplaybook}

\begin{lstlisting}[style=StyleTerminal,
  caption={Runbook --- confinement SSH compromis (Linux)},
  label={lst:runbook-ssh}]
# ─── Désactiver le compte compromis ──────────────────────────────────────────
# Sur le contrôleur de domaine (ou /etc/shadow si local)
usermod -L backup          # Verrouiller le compte Unix
# OU sur AD (PowerShell)
Disable-ADAccount -Identity "backup"

# ─── Bloquer l'IP source au firewall pfSense (CLI) ───────────────────────────
# Via l'API pfSense ou directement sur le shell pfSense
pfctl -t blocklist_ir -T add 203.0.113.42
pfctl -f /etc/pf.conf   # Recharger les règles

# ─── Isoler la machine dans un VLAN de quarantaine (switch manageable) ────────
# Changer le VLAN du port de srv-web-01 vers VLAN 99 (quarantaine)
# Via CLI switch (ex: Cisco IOS) :
# interface GigabitEthernet0/5
# switchport access vlan 99

# ─── Capturer la mémoire AVANT tout autre action ──────────────────────────────
ssh admin@10.0.1.10 "sudo avml /tmp/ram-$(hostname)-$(date +%Y%m%d-%H%M%S).lime"
scp admin@10.0.1.10:/tmp/ram-*.lime /mnt/evidence/
ssh admin@10.0.1.10 "sudo rm /tmp/ram-*.lime"

# ─── Documenter : toutes les commandes ci-dessus dans le log d'incident ───────
echo "$(date -u) | MINKA T.E. | Compte backup desactive, IP bloquee, VLAN quarantaine activé, RAM capturée" \
  >> /var/log/incident_INC-2026-0412-001.log
\end{lstlisting}

% ==============================================================================
\section{Confinement, éradication et rétablissement}
\label{sec:cir}
% ==============================================================================

\subsection{Confinement}

\begin{boxdef}
\textbf{Confinement~:} Mesures visant à limiter la propagation d'un incident
\textbf{sans altérer les preuves}. Le confinement est prioritaire sur
l'éradication~: on contient d'abord, on nettoie ensuite.
\end{boxdef}

Techniques de confinement~:
\begin{itemize}
  \item \textbf{Isolation réseau~:} VLAN de quarantaine, règle firewall,
        déconnexion câble (dernier recours).
  \item \textbf{Désactivation de compte~:} compte AD, compte local.
  \item \textbf{Blocage d'IP~:} au firewall, au niveau FAI si nécessaire.
  \item \textbf{Révocation de certificat / token~:} API keys, jetons OAuth.
  \item \textbf{Snapshot avant confinement~:} capturer l'état avant
        d'agir pour l'investigation.
\end{itemize}

\subsection{Éradication}

L'éradication supprime la cause de l'incident~:
\begin{itemize}
  \item Supprimer les fichiers malveillants identifiés.
  \item Désactiver / supprimer les mécanismes de persistance
        (Run keys, services, cron, authorized\_keys).
  \item Réinitialiser les credentials compromis (tous~: pas seulement
        le compte initialement compromis).
  \item Patcher la vulnérabilité exploitée.
  \item Scanner l'ensemble du parc pour détecter une propagation
        non identifiée.
\end{itemize}

\subsection{Rétablissement}

\begin{itemize}
  \item Restaurer les systèmes depuis des sauvegardes vérifiées
        (antérieures à la compromission).
  \item Remettre progressivement en production (pas en masse).
  \item Renforcer la surveillance~: règles SIEM spécifiques à l'attaque
        documentée, période de monitoring renforcée (30~jours).
  \item Valider le retour à la normale avec les équipes métier.
\end{itemize}

% ==============================================================================
\section{Rapport d'incident}
\label{sec:rapport-incident}
% ==============================================================================

\begin{boxdef}
\textbf{Rapport d'incident~:} Document structuré décrivant la nature
de l'incident, sa chronologie, son impact, les actions de réponse
entreprises et les recommandations. Deux versions~: technique (pour le
CSIRT, N3) et executive (pour la direction, non-technique).
\end{boxdef}

Structure type d'un rapport d'incident~:
\begin{enumerate}
  \item \textbf{Résumé exécutif} (1~page~: quoi, quand, impact, mesures).
  \item \textbf{Chronologie de l'incident} (timeline annotée MITRE).
  \item \textbf{Analyse technique} (vecteur initial, propagation, TTPs).
  \item \textbf{Preuves collectées} (liste avec hashes).
  \item \textbf{Impact} (technique~: systèmes affectés~; métier~:
        données, services, clients).
  \item \textbf{Mesures de réponse prises} (confinement, éradication,
        rétablissement).
  \item \textbf{Recommandations} (court terme~: remédiations~;
        long terme~: renforcement défensif).
  \item \textbf{Annexes} (logs bruts, captures, chaîne de custody).
\end{enumerate}

% ==============================================================================
\section{Mapping MITRE ATT\&CK}
\label{sec:mitre-ir}
% ==============================================================================

\begin{boxdef}
\textbf{MITRE ATT\&CK~:} Base de connaissance (knowledge base) des tactiques,
techniques et procédures (\termecle{TTPs}) utilisées par les attaquants dans
les environnements réels. 14~tactiques (colonnes), 200+~techniques (lignes),
organisées en matrices Enterprise, Mobile, ICS. Standard pour l'annotation
des rapports d'incident et l'ingénierie de détection.
\end{boxdef}

\begin{table}[H]
\centering\small
\caption{Mapping de l'incident LiZbiZ-SIS aux TTPs MITRE ATT\&CK}
\label{tab:mitre-mapping}
\rowcolors{2}{TableauPair}{TableauImpair}
\begin{tabular}{p{3.5cm}p{3.5cm}p{3.5cm}p{3.5cm}}
\toprule
\tableauHeader{Phase IR} & \tableauHeader{Action attaquant} &
\tableauHeader{TTP MITRE} & \tableauHeader{Détection SOC} \\
\midrule
Reconnaissance & Scan de ports & T1595 & pfSense log (SYN flood) \\
Accès initial & Brute force SSH & T1110.001 & Règle 100200 SIEM \\
Exécution & \texttt{wget} + script bash & T1059.004 & Sysmon EID~1 \\
C2 & Connexion HTTPS sortante & T1071.001 & Sysmon EID~3 + pfSense \\
Mouvement latéral & SMB vers ERP & T1021.002 & EID~4624 + EID~4776 \\
Collecte & Accès base de données & T1005 & Logs MySQL + proxy \\
Couverture & Effacement journal & T1070.001 & EID~1102 SIEM \\
\bottomrule
\end{tabular}
\end{table}

% ==============================================================================
\section*{Résumé}
% ==============================================================================

\begin{boxresume}
\begin{itemize}
  \item Cycle IR~: Préparation~$\to$~Détection~$\to$~Évaluation~$\to$
        Réponse~$\to$~RETEX (ISO~27035 / NIST~800-61).
  \item CSIRT~: équipe avec rôles définis (Manager, Analyste, Communication,
        Légal). Constituer avant l'incident.
  \item Playbook~: quoi faire (logique) ~; Runbook~: comment faire (commandes).
  \item Confinement d'abord~; preuves préservées~; éradication ensuite~;
        rétablissement progressif.
  \item Rapport d'incident~: résumé exécutif + chronologie + analyse
        technique + recommandations.
  \item Mapping MITRE ATT\&CK~: 14~tactiques, documenter chaque action
        attaquant avec son TTP.
\end{itemize}
\end{boxresume}

\begin{boxexo}
\textbf{Ex.~8.1 — Rédaction de playbook.}
Rédigez un playbook complet pour un incident de ransomware affectant
LiZbiZ-SIS~: 6~étapes, décisions à chaque étape, escalades, outils.

\medskip\textbf{Ex.~8.2 — Classification.}
À 02h30, une alerte Wazuh signale une modification de \texttt{/etc/passwd}
sur srv-web-01 par le compte \texttt{www-data}. Classifiez l'incident,
justifiez, et listez les 5~premières actions à mener.

\medskip\textbf{Ex.~8.3 — Rapport d'incident.}
Rédigez le résumé exécutif (1~page) de l'incident LiZbiZ-SIS du
12/04/2026 à destination du Directeur Général, sans jargon technique,
en couvrant~: nature de l'incident, impact, mesures prises,
coût estimé, recommandations.
\end{boxexo}
